% !TeX program = pdflatex
%
% A transformacao WebAssembly -> LLVM-IR no Map2Check
% Compilar no Overleaf com pdfLaTeX. Sem dependencias fora do TeX Live padrao.
%
\documentclass[11pt,a4paper]{article}

\usepackage[T1]{fontenc}
\usepackage[utf8]{inputenc}
\usepackage[brazil]{babel}
\usepackage{lmodern}
\usepackage{microtype}

% Caminhos longos em \texttt nao quebram; isto evita que estourem a margem
% sem precisar de \sloppy global.
\setlength{\emergencystretch}{3.5em}

\usepackage[margin=2.8cm]{geometry}
\usepackage{amsmath,amssymb,amsthm}
\usepackage{mathtools}
\usepackage{booktabs}
\usepackage{tabularx}
\usepackage{listings}
\usepackage{xcolor}
\usepackage{enumitem}
\usepackage[hidelinks]{hyperref}
\urlstyle{tt}   % \path{} para caminhos longos
\def\UrlBreaks{\do\/\do\.\do\-\do\_\do\{\do\}}

%---------------------------------------------------------------- ambientes
\theoremstyle{definition}
\newtheorem{hipotese}{Hipótese}
\newtheorem{proposicao}{Proposição}
\renewcommand{\qedsymbol}{$\blacksquare$}

%---------------------------------------------------------------- macros
\newcommand{\sem}[1]{[\![#1]\!]}          % colchetes semanticos, sem stmaryrd
\newcommand{\Safe}{\mathrm{Safe}}
\newcommand{\Reach}{\mathit{Reach}}
\newcommand{\wasm}{\mathrm{wasm}}
\newcommand{\lin}{\mathrm{lin}}
\newcommand{\code}[1]{\texttt{#1}}

% Regra de inferencia, sem dependencia externa: premissas sobre a barra,
% conclusao embaixo, nome da regra a direita.
\newcommand{\rulelabel}[1]{\textnormal{\textsc{\footnotesize #1}}}

%---------------------------------------------------------------- listings
\definecolor{cmtgray}{gray}{0.42}
\lstset{
  basicstyle=\ttfamily\small,
  commentstyle=\color{cmtgray}\itshape,
  keywordstyle=\bfseries,
  showstringspaces=false,
  columns=fullflexible,
  breaklines=true,
  frame=leftline,
  framesep=8pt,
  xleftmargin=10pt,
  captionpos=b
}

%---------------------------------------------------------------- titulo
\title{\textbf{A transformação WebAssembly $\rightarrow$ LLVM-IR no Map2Check}\\[0.4em]
       \large Formalização da cadeia de \emph{lifting}, e o que ela\\
       demonstravelmente preserva e perde}
\author{Projeto Map2Check}
\date{4 de setembro de 2026}

\begin{document}
\sloppy   % prefere espacamento folgado a texto invadindo a margem
\maketitle

\begin{abstract}
\noindent
Formalizamos a cadeia de transformações que leva um binário \code{.wasm} até o
programa instrumentado que o Map2Check analisa, decompondo-a em cinco funções e
enunciando explicitamente as duas hipóteses de correção sobre as quais ela
repousa. Dois resultados decorrem da formalização. A cadeia é \emph{sound} para
as dez classes de \emph{trap} do WebAssembly, o que já vale na implementação
atual e não vinha sendo reivindicado. E a cadeia é incompleta para segurança de
memória intra-\emph{sandbox} --- precisamente a classe de vulnerabilidade
declarada como alvo do projeto --- por uma razão estrutural: a memória linear
chega ao verificador como uma única alocação, de modo que a sobrescrita de um
objeto vizinho satisfaz tanto a verificação de limites do tradutor quanto o
predicado de segurança do próprio Map2Check. Damos o contraexemplo por
construção e derivamos, do enunciado, o requisito exato que uma modelagem de
memória linear teria de satisfazer.
\end{abstract}

\tableofcontents
\bigskip

%================================================================
\section{Objetivo e escopo}

Este documento descreve, passo a passo e com formalismo, a cadeia de
transformações que leva um binário \code{.wasm} até o programa instrumentado que
o Map2Check analisa.

O objetivo não é apenas documentar a implementação. É tornar \textbf{enunciável}
o que a cadeia preserva e o que ela perde --- porque essa distinção decide quais
resultados experimentais são atribuíveis ao Map2Check e quais não são. Duas
proposições saem dessa formalização, e ambas mudam o que se pode alegar sobre a
frente WASM.

A semântica de referência do WebAssembly é a de Haas et al., \emph{Bringing the
Web up to Speed with WebAssembly} (PLDI 2017), e a especificação normativa dela
derivada. O tradutor é o \code{wasm2c} do WABT 1.0.41. Todas as afirmações sobre
o código gerado foram verificadas na saída real do WABT na imagem
\path{ghcr.io/hbgit/map2check-dev}, não inferidas da documentação.

Referências de implementação: \path{modules/frontend/wasm_lifter.{hpp,cpp}},
\path{modules/frontend/caller.cpp} (função \code{generateWasmWrapperStatic}) e
\path{modules/frontend/map2check.cpp} (bloco \code{args.wasmMode}).

%================================================================
\section{Sumário dos resultados}

Para o leitor que precisa apenas da conclusão, antes do desenvolvimento.

\paragraph{Proposição 1 --- a cadeia é \emph{sound} para as dez classes de
\emph{trap} do WASM.}
O \code{wasm2c} compila a condição lateral da regra de redução de acesso à
memória em uma verificação dinâmica explícita, e essa verificação chama
\code{wasm\_rt\_trap}, que aborta. Abort é exatamente o sinal que o modo de
alcançabilidade do Map2Check já consome. Portanto acesso fora da memória linear,
divisão por zero, \emph{overflow} de conversão, \code{unreachable} e
\code{call\_indirect} inválido \textbf{são detectáveis hoje}, sem nenhum trabalho
adicional. Este resultado existe e não estava sendo alegado.

\paragraph{Proposição 2 --- a cadeia é incompleta para segurança de memória
intra-\emph{sandbox}.}
Esta é justamente a classe de vulnerabilidade declarada como alvo principal do
projeto. A memória linear chega ao verificador como \textbf{uma única alocação},
de modo que a sobrescrita de um objeto vizinho \emph{dentro} dela satisfaz tanto
a verificação de limites do WABT quanto o predicado de segurança do próprio
Map2Check. A demonstração é por construção e o contraexemplo tem quatro linhas de
C. A causa é estrutural, não um defeito de implementação.

\paragraph{Consequência imediata.}
A validação existente em \path{docs/reports/wasm-juliet-validation.csv} não separa
os dois mecanismos possíveis de detecção, e a Proposição~2 mostra que um deles é
impossível no caso intra-memória. Enquanto essa separação não for feita, os dados
são consistentes com a hipótese de que \textbf{o WABT encontrou os defeitos e o
Map2Check apenas reportou o abort resultante}.

%================================================================
\section{Notação}

Seja $\mathbb{B} = \{0,\dots,255\}$ o conjunto dos bytes. Escrevemos $[a,b)$ para
o intervalo inteiro semiaberto. Para uma função parcial $f$, $\operatorname{dom}(f)$
é seu domínio.

O tamanho de página do WebAssembly é a constante
\[
  P = 65\,536 .
\]

Os tipos de valor são
$\mathit{valtype} = \{\mathtt{i32}, \mathtt{i64}, \mathtt{f32}, \mathtt{f64}\}$,
e $\mathit{Val}$ denota o conjunto dos valores correspondentes.

%================================================================
\section{Passo 0 --- o objeto de entrada}

Um módulo WebAssembly é a tupla
\[
  M = (\mathit{types},\ \mathit{funcs},\ \mathit{tables},\ \mathit{mems},\
       \mathit{globals},\ \mathit{elem},\ \mathit{data},\ \mathit{start},\
       \mathit{imports},\ \mathit{exports}).
\]

Três componentes importam para o que segue.

\begin{itemize}[leftmargin=1.4em]
  \item $\mathit{mems}$, que na versão 1.0 do WebAssembly tem
        \textbf{cardinalidade no máximo 1}. Esta é a hipótese estrutural que
        sustenta todo o resto do documento: existe uma única memória linear por
        módulo.
  \item $\mathit{exports} : \mathit{Name} \rightharpoonup \mathit{ExternIdx}$,
        que dá as funções visíveis de fora do módulo.
  \item $\mathit{imports}$, que dá as funções fornecidas pelo \emph{host}. Em
        contexto de IoT, é por aqui que passam acesso a sensores, rede e banco de
        dados.
\end{itemize}

O \textbf{estado de execução} é o par $(s,F)$, com \emph{store}
\[
  s = (\mathit{funcs},\ \mathit{tables},\ \mathit{mems},\ \mathit{globals})
\]
e \emph{frame} $F$ contendo os locais e a instância corrente. Uma configuração é
$s; F; e^{*}$, onde $e^{*}$ é a pilha de instruções, e a semântica é dada por uma
relação de redução em passos pequenos
\[
  s; F; e^{*} \hookrightarrow s'; F'; e'^{*} .
\]
O termo distinguido $\mathtt{trap}$ é absorvente: nenhuma regra o reduz.

\subsection{A memória linear}

Uma instância de memória é
\[
  \mu = (d,\ \ell), \qquad d : [0,\ell) \to \mathbb{B}, \qquad \ell = n \cdot P
\]
para algum $n \in \mathbb{N}$.

Isto é: \textbf{um vetor plano de bytes, contíguo, endereçado por inteiros de 32
bits}, sem qualquer estrutura interna. Não há noção de objeto, de alocação, nem
de fronteira entre variáveis. Este fato é o centro de tudo que segue, e convém
insistir nele: a memória linear não é um espaço de endereçamento com objetos
dentro; \emph{é um único objeto}.

A regra de leitura, na forma que nos interessa, é

\begin{equation*}
  \frac{\;\; ea = i + \mathit{offset} \qquad ea + N/8 \le \ell \;\;}
       {\;\; s;\ F;\ (\mathtt{i32.const}\ i)\ (t.\mathtt{load}\ N\ \mathit{offset})
          \ \hookrightarrow\ s;\ F;\ (t.\mathtt{const}\ c) \;\;}
  \tag*{\rulelabel{load}}
\end{equation*}

\noindent
e, quando a condição lateral falha,

\begin{equation*}
  \frac{\;\; ea + N/8 > \ell \;\;}
       {\;\; s;\ F;\ (\mathtt{i32.const}\ i)\ (t.\mathtt{load}\ N\ \mathit{offset})
          \ \hookrightarrow\ s;\ F;\ \mathtt{trap} \;\;}
  \tag*{\rulelabel{load-trap}}
\end{equation*}

Convém guardar a forma da condição lateral, $ea + N/8 \le \ell$. Ela reaparece
literalmente, como código C, na seção seguinte.

%================================================================
\section{Passo 1 --- $\tau$: o tradutor \code{wasm2c}}

Definimos
\[
  \tau : \mathit{Wasm} \longrightarrow \mathit{C}
\]
como a função implementada pelo \code{wasm2c} do WABT 1.0.41, invocada em
\code{WasmLifter::runWasm2c} por

\begin{lstlisting}
wasm2c <entrada>.wasm -o <saida>.c
\end{lstlisting}

\noindent
produzindo o par de arquivos \code{.c} e \code{.h}.

\subsection{Representação do estado}

$\tau$ realiza o \emph{store} $s$ como uma \textbf{estrutura de instância} $I$. A
memória linear $\mu = (d,\ell)$ torna-se o registro \code{wasm\_rt\_memory\_t},
cujos campos relevantes são
\[
  \code{data} \mapsto d, \qquad \code{size} \mapsto \ell,
\]
com \code{data} obtido de uma \textbf{única} chamada de alocação, em
\code{wasm\_rt\_allocate\_memory}. Registre-se o quantificador: \emph{uma}
alocação, para \emph{toda} a memória do módulo. É desse fato que a Proposição~2
decorre.

Cada função $f$ do módulo torna-se uma função C
\[
  \tau(f) : \code{w2c\_}\langle \mathit{mod} \rangle \code{*}
            \times \overline{\mathit{Val}} \longrightarrow \mathit{Val},
\]
nomeada segundo o esquema \code{w2c\_<módulo>\_<função>}. O primeiro parâmetro é
o ponteiro de instância: é assim que o estado global do WASM se torna estado
explícito de primeira classe no C.

\subsection{A condição lateral torna-se código}

Eis o ponto técnico central. $\tau$ \textbf{compila a condição lateral da regra de
redução em uma verificação dinâmica explícita}. No C gerado, verificado na saída
real do WABT 1.0.41:

\begin{lstlisting}[language=C]
#define MEM_ADDR(mem, addr, n)  &((mem)->data[addr])

#define RANGE_CHECK(mem, offset, len)              \
  do {                                             \
    uint64_t res;                                  \
    if (UNLIKELY(add_overflow(offset, len, &res))) \
      TRAP(OOB);                                   \
    if (UNLIKELY(res > (mem)->size))               \
      TRAP(OOB);                                   \
  } while (0);

#define TRAP(x) (wasm_rt_trap(WASM_RT_TRAP_##x), 0)
\end{lstlisting}

Comparando com a regra \textsc{load}: \code{res > mem->size} é exatamente a
negação de $ea + N/8 \le \ell$, e \code{TRAP(OOB)} é exatamente $\mathtt{trap}$.
A tradução é literal.

Definindo então o predicado
\[
  \Safe_{\wasm}(\mu,\ ea,\ N) \;\equiv\; ea + N/8 \le \ell,
\]
o programa C avalia $\Safe_{\wasm}$ em tempo de execução antes de cada acesso,
desviando para \code{wasm\_rt\_trap} quando ele é falso.

O conjunto de \emph{traps} que o WABT distingue é
\[
  \mathcal{T} =
  \left\{
  \begin{array}{l}
    \mathtt{OOB},\quad \mathtt{INT\_OVERFLOW},\quad \mathtt{DIV\_BY\_ZERO},\quad
    \mathtt{INVALID\_CONVERSION},\\[2pt]
    \mathtt{UNREACHABLE},\quad \mathtt{CALL\_INDIRECT},\quad \mathtt{NULL\_REF},\\[2pt]
    \mathtt{UNCAUGHT\_EXCEPTION},\quad \mathtt{UNALIGNED},\quad \mathtt{EXHAUSTION}
  \end{array}
  \right\} .
\]

\subsection{Correção de \texorpdfstring{$\tau$}{tau}}

Seja $\approx$ a relação entre configurações WASM e estados C induzida pela
representação acima: igualdade byte a byte da memória, igualdade de valor para
globais e locais.

\begin{hipotese}[correção do \code{wasm2c}]
Para todo módulo $M$ bem tipado e toda execução, $\tau$ é uma simulação forte
módulo passos administrativos: se $s;F;e^{*} \approx \sigma$ e
$s;F;e^{*} \hookrightarrow s';F';e'^{*}$, então existe $\sigma'$ com
$\sigma \to^{+} \sigma'$ e $s';F';e'^{*} \approx \sigma'$; e
$e^{*} = \mathtt{trap}$ se e somente se $\sigma$ chama \code{wasm\_rt\_trap}.
\end{hipotese}

Isto é hipótese, \textbf{não teorema nosso}. Assumimos a correção do WABT como
assumimos a do Clang. Vale explicitá-la por duas razões: toda a \emph{soundness}
da cadeia repousa sobre ela, e ela é falsificável --- uma divergência do
\code{wasm2c} apareceria como discordância entre executar o \code{.wasm} num
\emph{runtime} conforme e executar o C gerado.

%================================================================
\section{Passo 2 --- $\kappa$: Clang para LLVM-IR}

\[
  \kappa : \mathit{C} \longrightarrow \mathit{LLVM}
\]
implementada em \code{WasmLifter::runClang} e \code{WasmLifter::runClangEmitLLVM}:

\begin{lstlisting}
clang-16 -c -emit-llvm -I<wasm-rt> -I<dir do .h>  <arq>.c -o <arq>.bc
clang-16 -S -emit-llvm -I<wasm-rt> -I<dir do .h>  <arq>.c -o <arq>.ll
\end{lstlisting}

A emissão do \code{.ll} textual não é redundante: ele é a entrada de $\epsilon$,
o passo seguinte.

\begin{hipotese}[correção do Clang]
$\kappa$ preserva a semântica de programas C \textbf{livres de comportamento
indefinido}.
\end{hipotese}

A ressalva importa mais aqui do que no uso normal do Map2Check. O C emitido pelo
\code{wasm2c} é gerado, não escrito, e emprega construções deliberadamente
próximas do limite da norma: aritmética de ponteiro sobre \code{data},
\code{setjmp}/\code{sigsetjmp} para o mecanismo de \emph{trap}, e o qualificador
\code{\_\_seg\_gs} quando o modo \emph{segue} está habilitado. Não auditamos essa
fronteira.

%================================================================
\section{Passo 3 --- $\epsilon$: extração do ponto de entrada}

\[
  \epsilon : \mathit{LLVM}_{\text{texto}} \longrightarrow
             \mathit{Name} \times \mathcal{P}(\mathit{Name})
\]
implementada em \code{WasmLifter::extractMetadata}. Ela percorre o \code{.ll}
linha a linha com a expressão regular

\begin{lstlisting}
define.*@w2c_([A-Za-z0-9_]+)\(
\end{lstlisting}

\noindent
e devolve o par $(\eta, E)$: o ponto de entrada $\eta$, reconhecido pelo sufixo
\code{\_0x5Fstart}, e o conjunto $E$ das demais funções exportadas.

Trata-se de \textbf{análise sintática sobre texto, não sobre o módulo LLVM}. Suas
condições de falha merecem enunciado, porque não são hipotéticas:

\begin{enumerate}[leftmargin=1.6em]
  \item Se nenhuma função casa com \code{\_0x5Fstart} --- o caso de um módulo
        \emph{reactor}, sem \code{\_start}, que é a forma usual de um
        \emph{handler} WASM de IoT --- o código adota o valor padrão
        \code{w2c\_\_start} e emite um aviso. Esse símbolo pode não existir.
  \item Por construção, $E$ e $\eta$ são disjuntos. Logo um módulo cujo único
        ponto de interesse são funções exportadas --- de novo, o caso IoT --- tem
        $\eta$ mal definido.
\end{enumerate}

Ou seja: $\epsilon$ é total, mas \textbf{não é correta} para a classe de módulos
que a proposta de pesquisa declara como alvo.

%================================================================
\section{Passo 4 --- $\omega$: o \emph{harness}}

O programa liftado não possui \code{main}. $\tau$ traduziu funções de módulo, e
nenhuma delas é ponto de entrada de um processo. Define-se então
\[
  \omega : \mathit{Name} \longrightarrow \mathit{LLVM}
\]
implementada em \code{Caller::generateWasmWrapperStatic}, que sintetiza
literalmente:

\begin{lstlisting}[language=C]
int main() {
    w2c_<mod> instance;
    wasm2c_<mod>_instantiate(&instance, NULL);
    <entrada>(&instance);
    wasm2c_<mod>_free(&instance);
    return 0;
}
\end{lstlisting}

O resultado é religado ao módulo liftado por
\[
  \Lambda_0 = \code{llvm-link}\bigl(\kappa(\tau(M)),\ \omega(\eta)\bigr).
\]

\subsection{O que $\omega$ quantifica: nada}

Formalmente, o módulo denota uma \textbf{família} de execuções, indexada pelos
argumentos das funções exportadas e pelo comportamento das importações:
\[
  \sem{M} : \overline{\mathit{Val}} \times \mathit{Host} \longrightarrow
            \mathit{Outcome}.
\]

O \emph{harness} atual não instancia essa família: ele a evita. A chamada
$\eta(\&I)$ é nulária --- o único argumento aceito é o ponteiro de instância,
porque $\eta$ é o \code{\_start}, cuja assinatura no WASM é
$[\,] \rightarrow [\,]$. Nenhuma valoração $\nu$ é escolhida porque não há onde
colocá-la.

A consequência é precisa. Escrevendo $\Reach(p)$ para o conjunto de estados
alcançáveis de um programa $p$:
\[
  \Reach(\Lambda_0)
  \;=\; \bigcup_{h \in \mathit{Host}} \Reach\bigl(\sem{M}(\bot,\ h)\bigr)
  \;\subseteq\; \bigcup_{\nu,\ h} \Reach\bigl(\sem{M}(\nu,\ h)\bigr),
\]
e a inclusão é \textbf{estrita} sempre que o módulo possui alguma função exportada
com aridade positiva. O símbolo $\bot$ não denota um valor: denota a ausência da
chamada.

Duas leituras decorrem disso, e ambas importam.

\paragraph{Para um módulo com \code{\_start}.}
Um programa compilado para WASI, que é o formato dos casos do Juliet na validação
existente. A entrada chega por \code{argv} ou por \code{stdin} \emph{dentro} do
\code{\_start}, e o \code{NonDetPass} pode amarrá-la ali. A cadeia funciona. Este
é o caso efetivamente medido até hoje.

\paragraph{Para um módulo \emph{reactor}.}
O \emph{handler} de IoT que a proposta declara como alvo, sem \code{\_start},
cujo comportamento inteiro reside nas funções exportadas em $E$. O \emph{wrapper}
não chama nada de interesse. Combinado com a falha de $\epsilon$ descrita na seção
anterior, o resultado é que \textbf{o alvo declarado da proposta não é alcançável
pela implementação atual}. E não por um motivo profundo: falta gerar uma chamada
por função exportada, com \code{\_\_VERIFIER\_nondet\_<t>()} em cada parâmetro.

O tratamento das importações $h$ permanece como buraco restante. Sem um modelo de
\emph{host}, as funções importadas são símbolos externos, e o que a análise assume
sobre elas é o comportamento padrão de quem estiver executando, não uma escolha
nossa. A fase de análise de \emph{taint} prevista na proposta --- rastrear dado
controlado até uma função importada --- vive inteiramente nesse buraco.

%================================================================
\section{Passo 5 --- $\pi$: a instrumentação do Map2Check}

Deste ponto em diante o pipeline é o comum: o bitcode religado é tratado como
qualquer entrada em LLVM-IR.

$\pi$ é a transformação de programa realizada pelos passes LLVM. Ela é um par
$(\pi_{\text{code}},\ \Sigma)$: instrumentação e \textbf{estado sombra}. No modo
\code{--memtrack}, $\Sigma$ mantém o conjunto das alocações vivas
\[
  \Sigma \subseteq
  \bigl\{ (\mathit{base},\ \mathit{size},\ \mathit{alive}) \bigr\},
\]
e $\pi_{\text{code}}$ insere, antes de cada acesso ao par $(p,n)$, a verificação
\[
  \Safe_{\mathrm{C}}(\Sigma,\ p,\ n) \;\equiv\;
  \exists\, (b,z,\mathit{alive}) \in \Sigma :\;
  \mathit{alive} \ \wedge\ b \le p \ \wedge\ p + n \le b + z .
\]

Compare este predicado com $\Safe_{\wasm}$. \textbf{As duas verificações não falam
da mesma coisa}, e as duas seções seguintes são inteiramente sobre a distância
entre elas.

A cadeia completa é
\[
  \boxed{\;
  \Lambda \;=\; \pi \;\circ\; \operatorname{link}\bigl(\omega \circ \epsilon,\ \cdot\bigr)
          \;\circ\; \kappa \;\circ\; \tau
  \;}
\]
seguida do motor de geração de entradas --- KLEE ou LibFuzzer --- sobre
$\Lambda(M)$.

%================================================================
\section{Proposição 1: o que a cadeia preserva}

\begin{proposicao}[\emph{soundness} para alcançabilidade de \emph{trap}]
\label{prop:sound}
Sob as Hipóteses 1 e 2, se a execução de $\Lambda(M)$ alcança
\code{wasm\_rt\_trap}, então a execução correspondente de $M$ reduz a
$\mathtt{trap}$.
\end{proposicao}

\begin{proof}[Demonstração (esboço)]
Pela Hipótese 2, a execução de $\Lambda_0$ corresponde à do C gerado. Pela
Hipótese 1, alcançar \code{wasm\_rt\_trap} no C corresponde a alcançar
$\mathtt{trap}$ no WASM. A transformação $\pi$ apenas acrescenta código operando
sobre um estado sombra disjunto do estado do programa, logo preserva
alcançabilidade.
\end{proof}

O resultado é real e é útil. \textbf{As dez classes de \emph{trap} de
$\mathcal{T}$ são detectáveis pela cadeia atual, sem trabalho adicional}: acesso
fora da memória linear, divisão por zero, \emph{overflow} de conversão, conversão
inválida a partir de NaN, \code{unreachable}, \code{call\_indirect} inválido,
referência nula, acesso atômico desalinhado e esgotamento de pilha. O que
\code{wasm\_rt\_trap} faz é abortar, e abort é precisamente o sinal que o modo de
alcançabilidade do Map2Check já consome.

%================================================================
\section{Proposição 2: o que a cadeia não preserva}

\begin{proposicao}[incompletude para segurança de memória]
\label{prop:incompleta}
Existem um módulo $M$ e uma execução tais que $M$ \textbf{sobrescreve um objeto
que não é o alvo pretendido} --- uma violação genuína de CWE-787 --- e, ainda
assim, $\Lambda(M)$ não viola $\Safe_{\mathrm{C}}$ em nenhum acesso, nem alcança
\code{wasm\_rt\_trap}.
\end{proposicao}

\begin{proof}[Demonstração (por construção)]
Tome o módulo compilado a partir de

\begin{lstlisting}[language=C]
int  buf[4];
char token[16];
int f(int i) { buf[i] = 7; return buf[0]; }
\end{lstlisting}

\noindent
com $i = 4$, e com $\mathit{token}$ posicionado imediatamente após $\mathit{buf}$
na memória linear. Sejam $a_{\mathit{buf}}$ e $a_{\mathit{token}}$ os respectivos
deslocamentos.

\medskip
\noindent\textbf{No WASM.}
O acesso tem endereço efetivo $ea = a_{\mathit{buf}} + 16 = a_{\mathit{token}}$.
Como $ea + 4 \le \ell$ --- a memória linear tem ao menos uma página, logo
$\ell \ge 65\,536$, e os objetos estão no início dela --- a condição lateral da
regra \textsc{load} \emph{é satisfeita}. A redução não produz $\mathtt{trap}$:
produz uma escrita bem definida. O \emph{token} foi sobrescrito e a execução
prossegue normalmente.

\medskip
\noindent\textbf{No C gerado.}
\code{RANGE\_CHECK} testa \code{res > mem->size}, isto é $ea + 4 > \ell$, que é
falso. Nenhum \emph{trap}.

\medskip
\noindent\textbf{Sob a instrumentação.}
Como visto, \code{data} é uma \emph{única} alocação cobrindo $[0,\ell)$. Logo
$\Sigma$ contém o registro $(\code{data},\ \ell,\ \mathrm{true})$ e, para
$p = \code{data} + ea$ e $n = 4$, valem $\code{data} \le p$ e
$p + 4 \le \code{data} + \ell$. Portanto $\Safe_{\mathrm{C}}(\Sigma, p, 4)$ é
verdadeiro e nenhuma violação é reportada.
\end{proof}

O ponto não é uma limitação do modo \code{memtrack}. É que \textbf{$\tau$ apaga a
informação necessária}: a partição da memória em objetos existe no fonte C
original, sobrevive até o compilador para WASM, e é descartada na emissão do
módulo, porque o WebAssembly não tem onde guardá-la. O que chega ao Map2Check é um
vetor de bytes com uma única fronteira --- a externa.

Note ainda o \textbf{quantificador invertido}: quanto \emph{maior} a memória
linear, \emph{menos} detectável a violação, porque mais distante fica a única
fronteira que restou. Um módulo com uma única página tem $\ell = 65\,536$; um
\emph{overflow} de dezesseis bytes está a sessenta e cinco mil bytes de se tornar
visível.

%================================================================
\section{Consequência para a validação existente}

O resultado acima muda como se deve ler
\path{docs/reports/wasm-juliet-validation.csv}. Todos os casos ali registrados
retornam \code{FALSE}, isto é, defeito encontrado. Mas as Proposições
\ref{prop:sound} e \ref{prop:incompleta}, tomadas juntas, dizem que há
\textbf{dois mecanismos possíveis} para esse veredito:

\begin{enumerate}[label=(\alph*),leftmargin=1.8em]
  \item o \emph{overflow} foi grande o bastante para sair de $[0,\ell)$, o
        \code{RANGE\_CHECK} do WABT disparou, e o que foi detectado foi a
        verificação do WABT; ou
  \item a instrumentação do Map2Check detectou algo por conta própria.
\end{enumerate}

A planilha \textbf{não distingue os dois casos}, e a Proposição
\ref{prop:incompleta} estabelece que (b) é impossível para violação
intra-memória. Enquanto essa separação não for feita, a validação WASM não
sustenta a alegação de que \emph{o Map2Check} encontrou os defeitos: ela é
igualmente consistente com a hipótese de que o WABT os encontrou e o Map2Check
apenas reportou o abort resultante.

Separar é barato: basta instrumentar \code{wasm\_rt\_trap} e registrar qual dos
dois caminhos produziu o abort. É a primeira medição que a frente WASM deveria
fazer, e ela precede qualquer alegação nova.

%================================================================
\section{O que a modelagem de memória linear tem de acrescentar}

A proposta de pesquisa prevê uma fase de ``modelagem da memória linear''. Nos
termos deste documento, o que falta é exatamente um objeto: um \textbf{mapa de
objetos}
\[
  O : [0,\ell) \rightharpoonup \mathit{ObjID},
\]
com o predicado refinado
\[
  \Safe_{\lin}(O,\ ea,\ N) \;\equiv\;
  \Safe_{\wasm}(\mu,\ ea,\ N)
  \ \wedge\
  \exists\, o \in \mathit{ObjID} :\
  \forall k \in [\,ea,\ ea + N/8\,) :\ O(k) = o .
\]

Em palavras: além de estar dentro da memória, o acesso inteiro precisa cair
\textbf{dentro de um único objeto}. Com $\Safe_{\lin}$ no lugar de
$\Safe_{\mathrm{C}}$, o contraexemplo da Proposição \ref{prop:incompleta} passa a
ser detectado, pois $O(a_{\mathit{buf}} + 16) = \mathit{token} \ne \mathit{buf}$.

Isso torna preciso o requisito sobre o \emph{lifter}. O mapa $O$ \textbf{não é
recuperável do \code{.wasm}}, porque foi apagado antes. As três fontes possíveis
são:

\begin{enumerate}[leftmargin=1.6em]
  \item \textbf{Seções de nome e informação DWARF} do módulo, quando presentes ---
        o Clang com \code{-g} para o alvo \code{wasm32} as emite. Recupera objetos
        estáticos e o \emph{layout} de quadro de pilha.
  \item \textbf{Reconstrução do alocador}: reconhecer as funções \code{malloc} e
        \code{free} do módulo, quando ele carrega uma libc, e delas derivar $O$
        dinamicamente.
  \item \textbf{Anotação no momento do \emph{lifting}}, que é o que a proposta
        chama de ``\emph{lifter} ciente de segurança'': um $\tau'$ que emita,
        junto do C, os metadados de fronteira.
\end{enumerate}

As três são trabalho de $\tau$, ou de um passo novo entre $\tau$ e $\kappa$.
\textbf{Nenhuma delas é trabalho do backend de verificação.} Isso confirma a
afirmação de viabilidade da proposta --- não é preciso reescrever o motor --- e ao
mesmo tempo localiza o custo real onde ele está.

%================================================================
\section{Resumo da cadeia}

\begin{table}[h]
\centering
\small
\renewcommand{\arraystretch}{1.3}
\begin{tabularx}{\textwidth}{@{}c l l X X@{}}
\toprule
\textbf{Passo} & \textbf{Função} & \textbf{Implementação} &
\textbf{Preserva} & \textbf{Perde} \\
\midrule
0 & $M$ & entrada \code{.wasm} & --- &
  o \emph{layout} de objetos, já perdido na compilação \emph{para} WASM \\
1 & $\tau$ & \code{wasm2c}, WABT 1.0.41 &
  semântica, \emph{traps}, bytes da memória & nada além do que já faltava \\
2 & $\kappa$ & \code{clang-16 -emit-llvm} &
  semântica C, módulo comportamento indefinido & --- \\
3 & $\epsilon$ & regex sobre o \code{.ll} &
  ponto de entrada, quando há \code{\_start} & módulos \emph{reactor}, o caso IoT \\
4 & $\omega$ & \emph{wrapper} e \code{llvm-link} &
  a execução de \code{\_start} &
  toda função exportada de aridade positiva, e o \emph{host} \\
5 & $\pi$ & passes do Map2Check & alcançabilidade & --- \\
\bottomrule
\end{tabularx}
\end{table}

%================================================================
\section{Recomendações, em ordem}

\begin{enumerate}[leftmargin=1.6em]
  \item \textbf{Separar (a) de (b).} Instrumentar \code{wasm\_rt\_trap} para
        registrar qual caminho produziu cada abort, e reprocessar a validação
        Juliet-WASM existente. Custo baixo, e nenhuma alegação sobre a frente WASM
        se sustenta antes disso.
  \item \textbf{Enunciar a Proposição \ref{prop:sound} como resultado.} A
        \emph{soundness} para as dez classes de \emph{trap} já existe na
        implementação atual e não está sendo reivindicada. É um resultado
        publicável por si, e independente da fase de modelagem de memória.
  \item \textbf{Corrigir $\epsilon$ e $\omega$ para módulos \emph{reactor}.} Gerar
        uma chamada por função exportada, com
        \code{\_\_VERIFIER\_nondet\_<t>()} em cada parâmetro. Sem isso, o alvo
        declarado da proposta --- o \emph{handler} de IoT --- permanece fora do
        alcance da ferramenta, e por um motivo puramente mecânico.
  \item \textbf{Construir o mapa de objetos $O$.} É o que dá conteúdo à alegação
        de detecção de corrupção de memória intra-\emph{sandbox}, e é a única das
        quatro que constitui pesquisa em vez de engenharia.
\end{enumerate}

\end{document}
